\section{Preliminaries}\label{sec:prelim}

\paragraph{Point sets and the range-counting oracle.}
The input is a set $P$ of $n\ge 2$ distinct points in the integer grid $[\Delta]^2$, where
$[\Delta]=\{0,1,\dots,\Delta-1\}$ and $\Delta=O(n)$; thus $\log\Delta=O(\log n)$. (For $n\le 1$ the MST
weight is $0$.) Distinctness gives $W=\wt(\MST(P))\ge n-1$, since every MST edge has length at least one. We access $P$ only
through an \emph{orthogonal range-counting oracle}: a query is an axis-aligned rectangle
$R=[a_1,b_1)\times[a_2,b_2)$, and the oracle returns $\rcount{R}=\abs{P\cap R}$. The cost of an algorithm
is the number of oracle queries it makes; all other computation is free and charged separately only in
passing (the algorithms here run in time near-linear in their query count on a standard word RAM with
$O(\log n)$-bit words). Two derived primitives are used throughout and each costs $O(\log\Delta)$
queries: an \emph{emptiness} test ($\rcount{R}\overset{?}{=}0$) and a \emph{uniform point sample} from
$P\cap R$ (binary search on nested rectangles; see \cref{sec:prelim-sample}). We write $\Otil(\cdot)$ to
hide factors polylogarithmic in $n$, $\Delta$, and the accuracy parameters.

\paragraph{The estimation problem.}
For $p,q\in\R^2$ let $\dist(p,q)$ denote Euclidean ($\ell_2$) distance, and let $\MST(S)$ be a Euclidean
minimum spanning tree of a finite point set $S$, with weight $\wt(\MST(S))=\sum_{e\in\MST(S)}\abs{e}$.
We study the query complexity of estimating $W\eqdef\wt(\MST(P))$.

\begin{problem}\label{prob:main}
Given $\eps\in(0,1)$ and oracle access to $P\subseteq[\Delta]^2$ with $\abs{P}=n$, output $\widehat W$
with $(1-\eps)W\le\widehat W\le(1+\eps)W$ with probability at least $2/3$, using as few range-counting
queries as possible.
\end{problem}

The success probability $2/3$ is boosted to $1-\delta$ by $O(\log(1/\delta))$ independent repetitions and
a median. Driemel, Monemizadeh, Oh, Staals, and Woodruff~\cite{DMOSW25} give an $\Otil(\sqrt n)$-query
algorithm for \cref{prob:main} and an $\Omega(n^{1/3})$-query lower bound, leaving a polynomial gap. We
close it from above.

The range-counting oracle is the weakest of the geometric oracles used for this problem and the one a
spatial index supports natively: it returns a cardinality, not a witness, so the algorithm never learns
the identity or location of any individual point except through the $O(\log\Delta)$-query sampling
primitive of \cref{sec:prelim-sample}. It is in particular weaker than the range-emptiness-with-nearest-
neighbor model of~\cite{Czumaj05}, which can locate points. Every quantity we estimate is ultimately a
function of the counts $\rcount R$ over a family of aligned rectangles, and the cost is the number of
distinct such queries.

\paragraph{The cluster-count integral.}
For a finite $S\subset\R^2$ and a threshold $t\ge 0$, let $G_S(t)$ be the graph on $S$ with an edge
between every pair of points at distance at most $t$, and let $c_S(t)$ be its number of connected
components. The components of $G_S(t)$ are exactly the \emph{single-linkage clusters} of $S$ at scale
$t$, and $c_S$ is a nonincreasing step function with $c_S(0)=\abs{S}$ (for $S$ in general position) and
$c_S(t)=1$ for $t\ge\diam(S)$. The Euclidean MST weight is the area under this curve:
\begin{equation}\label{eq:cluster-integral}
  \wt(\MST(S))=\int_0^\infty\bigl(c_S(t)-1\bigr)\,dt .
\end{equation}
Equation~\eqref{eq:cluster-integral} is the Euclidean instance of the Chazelle--Rubinfeld--Trevisan
identity~\cite{CRT05}, and we recall its one-line proof since the same threshold-graph view recurs
throughout. The key fact is
\begin{equation}\label{eq:components-edges}
  c_S(t)-1=\abs{\set{e\in\MST(S):\abs{e}>t}},\qquad\text{hence}\qquad c_S(t)-1\le \frac{\wt(\MST(S))}{t}.
\end{equation}
To see the left identity, run Kruskal's algorithm on $S$ in increasing edge length: it builds $\MST(S)$,
and at the moment it considers an edge $e$ of length $\abs e$, the components present are exactly the
clusters of $S$ at threshold just below $\abs e$, so adding the MST edges of length $>t$ is what reduces
the cluster count from $c_S(t)$ down to $1$. Each such edge lowers the count by one (it joins two distinct
clusters, by the cut property), so the number of MST edges longer than $t$ is $c_S(t)-1$. Integrating this
identity over $t\in[0,\infty)$ and exchanging sum and integral,
\[
  \int_0^\infty(c_S(t)-1)\,dt=\int_0^\infty\abs{\set{e\in\MST(S):\abs e>t}}\,dt
   =\sum_{e\in\MST(S)}\int_0^{\abs e}dt=\sum_{e\in\MST(S)}\abs e=\wt(\MST(S)),
\]
which is \eqref{eq:cluster-integral}. The right inequality of \eqref{eq:components-edges} is Markov on the
edge lengths ($t\cdot\abs{\{e:\abs e>t\}}\le\sum_e\abs e$) and is used repeatedly to bound component counts
at large thresholds. Estimating $W$ thus reduces to estimating the cluster-count curve $c_P(\cdot)$, which
is what the oracle must support.

\paragraph{Spanners and the death-time view.}
There are two ways to read the cluster-count curve, and we use both. The support estimator of
\cref{sec:support} reads it by sampling \emph{space}; for the bulk of $P$ we instead read it by sampling
\emph{vertices} of a sparse graph and measuring how far each must travel to merge with someone senior.
Concretely we sample from $P$ and explore a locally navigable $(1+\rho)$-spanner $H$ of the sampled
neighborhood, following~\cite{DMOSW25}. For a connected weighted graph $H$ on vertex set $V$ with a fixed
root $r$, assign independent uniform ranks to $V\setminus\{r\}$ and define the \emph{death time} $\tau(v)$
of a vertex as the bottleneck distance from $v$ to the root or to any lower-ranked vertex. The multiset
$\set{\tau(v):v\in V}$ equals $\set{0}\cup\set{\wt(e):e\in\MST(H)}$ (\cref{sec:dt}), so the clipped
statistics $\min\{\tau(v),L\}$ recover the truncated MST weight $\sum_{e}\min\{\abs{e},L\}$, and a uniform
sample of vertices estimates that truncated weight without ever forming the MST. The reason a spanner
suffices is that the threshold graphs of $H$ bracket those of the points: because $H$ is a Euclidean
$(1+\rho)$-spanner, $c_S(t)\le c_H(t)\le c_S(t/(1+\rho))$, so its death times approximate the point-set
MST within a $(1+\rho)$ factor. The catch is that we cannot build $H$ explicitly under a counting oracle;
we realize it implicitly with the well-separated pair decomposition (WSPD)~\cite{CK95,HarPeled11},
navigating one vertex's spanner neighborhood at a time. The full implementation is \cref{sec:wspd}.

We recall the two WSPD facts the implementation rests on. For a quadtree on $[\Delta]^2$ and a separation
parameter $1/\rho$, a pair of cells $(c,c')$ is \emph{well separated} when their centers are at distance at
least $\max\{r(c),r(c')\}/\rho$, where $r(\cdot)$ is the cell side; the Callahan--Kosaraju recursion
produces a family $\mathcal W$ of well-separated pairs such that every pair of points is separated by
exactly one pair of $\mathcal W$.

\begin{lemma}[WSPD properties \cite{CK95,HarPeled11}]\label{lem:wspd-prelim}
For separation $1/\rho$: (i) every $(c,c')\in\mathcal W$ has $r(c')/2\le r(c)\le 2r(c')$; (ii) every point
of $[\Delta]^2$ lies in $O(\rho^{-2}\log\Delta)$ pairs of $\mathcal W$; and (iii) connecting a fixed
representative of each cell across every pair, with the edge weight set to the representatives' Euclidean
distance, yields a $(1+O(\rho))$-spanner of the point set whose MST weight is at most $(1+O(\rho))$ times
the Euclidean MST weight.
\end{lemma}

Facts (i)--(iii) are standard~\cite{HarPeled11}; \cref{sec:wspd} turns them into a range-count
implementation, the only nonstandard point being that a counting oracle reveals neither the points nor the
WSPD, so each representative, each pair-membership test, and each spanner edge must be recovered by
queries.

\paragraph{Notation.}
Throughout, $K$ denotes the number of nonempty cells of a grid imposed on $P$ (the size of a regularized
\emph{support}), $Q$ the set of cell centers (the support point set), and $\Wq\eqdef\wt(\MST(Q))$. We
write $H_m=\sum_{i=1}^m 1/i=O(\log m)$ for the harmonic number, and reserve $\eta,\beta,\xi,\rho,\sigma$
for accuracy parameters, all in $(0,1)$. \Cref{tab:notation} collects the recurring symbols.

\begin{table}[t]
\centering
\small
\begin{tabular}{ll}
\toprule
symbol & meaning\\
\midrule
$P\subseteq[\Delta]^2$, $n=\abs P$ & input point set; $\Delta=O(n)$\\
$\rcount{R}$ & oracle answer $\abs{P\cap R}$ on rectangle $R$\\
$W=\wt(\MST(P))$ & quantity to estimate\\
$c_S(t)$ & \# components of the distance-$t$ threshold graph on $S$\\
$\mathcal G_h$, $h$ & aligned grid and its cell side\\
$Q$, $K=\abs Q$ & cell-center support and its size\\
$\Wq=\wt(\MST(Q))$ & support MST weight\\
$b$, $\delta$ & \# active-cover blocks and their side ($b=\Theta(\sqrt K)$)\\
$\Gamma_r$, $\cG(r)$ & scale graph at threshold $r$ and its \# components\\
$\tau(v)$ & death time of vertex $v$\\
$A_L,\,B_L$ & $L$-clipped and $L$-tail MST weight\\
\bottomrule
\end{tabular}
\caption{Recurring notation.}
\label{tab:notation}
\end{table}

\begin{table}[t]
\centering
\small
\begin{tabular}{lll}
\toprule
parameter & meaning & setting\\
\midrule
$\eps$ & target relative accuracy & input\\
$\zeta$ & per-call normalized accuracy & $\eps/32$ (final call)\\
$\xi$ & internal accuracy of one call & $\zeta/C_*$\\
$\rho$ & spanner stretch $1+\rho$ & $\Theta(\xi)$\\
$\eta$ & relative accuracy of \cref{lem:support} & $\Theta(\xi)$\\
$\beta$ & additive-call accuracy (support) & $\Theta(\eta)$\\
$\sigma$ & scale-graph fineness & $\beta/32$\\
$K_0$ & target support size & $\ceil{n^{2/3}}$\\
$L$ & clip level & $G/K_0$\\
\bottomrule
\end{tabular}
\caption{Accuracy parameters and their dependencies. Each is a fixed $\Theta(\cdot)$ multiple of the one
above it, so all are $\Theta(\eps)$ and every cost is $\Otil_\eps(\cdot)$; $C_*$ is the absolute constant
of the assembly bound \eqref{eq:guess-accuracy}.}
\label{tab:params}
\end{table}

\paragraph{Probabilistic tools.}
We use two standard facts repeatedly. The first is the multiplicative Chernoff bound: if $Y$ is a sum of
independent $[0,1]$ variables with mean $\mu$, then $\Pr[\abs{Y-\mu}\ge\gamma\mu]\le 2e^{-\gamma^2\mu/3}$
for $\gamma\in(0,1)$. The second is median-of-means amplification: to estimate a quantity $q$ to additive
error $\alpha$ with failure probability $\delta$ from an unbiased estimator of variance $V$, average
$O(V/\alpha^2)$ samples into one block (so each block is within $\alpha$ with probability $\ge 2/3$ by
Chebyshev) and take the median of $O(\log(1/\delta))$ blocks; the median is within $\alpha$ with
probability $1-\delta$. We invoke this both for the leader estimator (\cref{lem:leader}, with $V\le MU$)
and for the death-time averages (\cref{sec:dt-sampler}). When several estimates must hold at once we set
each block's $\delta$ to a $1/\poly$ fraction and union-bound, absorbing the $O(\log)$ overhead into
$\Otil(\cdot)$.
